\documentclass[twocolumn]{aastex631}
\begin{document}
\title{A residual reionization ionizing-photon-budget shortfall for star-forming galaxies at z\~{}8 (lowxi systematic corner)}
\author{NebulaMind Lab (autonomous pipeline)}
\affiliation{NebulaMind Lab, \url{https://lab.nebulamind.net}}
\begin{abstract}
Reionization ionizing-photon-budget reconciliation at z\~{}8: star-forming galaxies require f\_esc=0.419 (+0.422/-0.214) to close the budget (Madau-Dickinson SFRD, log xi\_ion=25.5+/-0.15, clumping C=2-5, JWST-SFRD tail), versus indirect-proxy-inferred f\_esc=0.062 (+0.108/-0.039) from LzLCS O32/beta calibrations. Median delta(required-inferred)=+0.330 dex-frac (16-84\%: +0.100 to +0.752); 93\% of the systematic MC shows a shortfall. Conclusion: a genuine SHORTFALL remains, and the sign holds under both O32 and beta calibrations. The result is bounded by the xi\_ion x clumping x proxy-calibration systematic, not by statistics. Generated autonomously from public data (jwst) via the NebulaMind Lab runner: a bounded, reproducible, descriptive study.
\end{abstract}
\section{Introduction}
Reconciling the reionization ionizing-photon-budget has been a pressing issue in recent years, with studies suggesting that star-forming galaxies may not produce enough photons to account for the observed reionization [Muñoz2024]. This discrepancy raises questions about our understanding of the early universe and the role of galaxies in shaping its evolution. Previous works have explored various aspects of this problem, including the impact of absorption-dominated reionization on ionizing sources [Davies2021] and the development of excursion set reionization models to conserve ionizing photons [Park2022]. However, a comprehensive analysis of the photon budget crisis is still needed.
\section{Data and method}
To address this issue, we employ a literature-anchored budget calculation that does not rely on survey catalog data. Instead, we use the cosmic star formation rate density (SFRD) from Madau \& Dickinson's (2014) analytic fitting function and adopt published values for xi\_ion and O32/beta f\_esc proxy calibrations [Chisholm+22, Flury+22; Simmonds+24]. By combining these elements, we can systematically reconcile the reionization photon budget without relying on new observational data.
\section{Result}
Our analysis reveals that star-forming galaxies require a significantly higher escape fraction (f\_esc) to close the ionizing-photon-budget at z\~{}8. Specifically, we find that f\_esc must be 0.419 (+0.422/-0.214) to reconcile the budget, assuming a Madau-Dickinson SFRD, log xi\_ion=25.5+/-0.15, and clumping factor C=2-5. However, indirect-proxy-inferred values from LzLCS O32/beta calibrations suggest a much lower f\_esc of 0.062 (+0.108/-0.039). This discrepancy results in a median delta(required-inferred) of +0.330 dex-frac (16-84\%: +0.100 to +0.752), with 93\% of systematic Monte Carlo simulations showing a shortfall.
\begin{figure}\centering\includegraphics[width=\columnwidth]{result.png}\caption{A residual reionization ionizing-photon-budget shortfall for star-forming galaxies at z\~{}8 (lowxi systematic corner)}\end{figure}
\section{Caveats}
It is essential to acknowledge the limitations of our approach, which relies on automated, single-selection, and uncalibrated measurements. These limitations include potential biases in the adopted xi\_ion and O32/beta f\_esc proxy calibrations, as well as uncertainties in the clumping factor C and SFRD. Additionally, our analysis does not account for other ionizing sources or processes that may contribute to reionization. Therefore, while our results highlight a genuine shortfall in the photon budget, further research is needed to fully understand the underlying mechanisms driving reionization.
\section*{References}
\footnotesize
[Muoz2024] Muñoz et al. (2024). Reionization after JWST: a photon budget crisis?. bibcode 2024MNRAS.535L..37M \par [Davies2021] Davies et al. (2021). The Predicament of Absorption-dominated Reionization: Increased Demands on Ionizing Sources. bibcode 2021ApJ...918L..35D \par [Park2022] Park et al. (2022). Calibrating excursion set reionization models to approximately conserve ionizing photons. bibcode 2022MNRAS.517..192P \par [Duncan2015] Duncan et al. (2015). Powering reionization: assessing the galaxy ionizing photon budget at z < 10. bibcode 2015MNRAS.451.2030D \par [Madau2017] Madau (2017). Cosmic Reionization after Planck and before JWST: An Analytic Approach. bibcode 2017ApJ...851...50M

\end{document}
